% META: {"id": "TNSI-PROJET-AM-EXTRAIT", "chapitre": "TNSI-PROJET", "type_objet": "amenagee", "capacites_codes": ["C1", "C2"], "derive_de": ["TNSI-PROJET-ANNUEL"], "profil_amenagement": ["SEQUENCAGE", "ALLEGEMENT_ENONCE", "REPONSE_FERMEE", "AERATION", "AMORCE_FOURNIE", "RETRAIT_DOUBLE_TACHE"], "mode_creation": "adaptation", "sources_inspiration": [], "status": "generated"}
\section*{Version aménagée --- Extrait Projet annuel}
\textit{Consignes séquencées, mise en page aérée, énoncés allégés.}

\medskip
\textit{Cette fiche ne remplace pas le dossier de projet : elle en séquence la
préparation. Les exigences évaluées restent celles de la grille critériée.}

\bigskip

\subsection*{Partie A --- Poser le besoin (C1)}

\textbf{Étape 1.} Écris en \textbf{une seule phrase} ce que ton programme permettra
de faire.

\medskip
« Mon programme permet à \dotfill\ de \dotfill\ »

\medskip
\textit{Contrôle :} si ta phrase ne tient pas sur une ligne, le projet n'est
pas encore assez précis.

\bigskip

\textbf{Étape 2.} Coche ce que tu remettras à la fin.

\bigskip

\begin{center}
\begin{tabular}{|l|c|}
\hline
\textbf{Livrable} & \textbf{Prévu ?} \\
\hline
Le programme & oui \\
\hline
Un jeu de tests & \dotfill \\
\hline
Une documentation & \dotfill \\
\hline
Une présentation orale & \dotfill \\
\hline
\end{tabular}
\end{center}

\bigskip
\bigskip

\subsection*{Partie B --- Découper en jalons (C2)}

\textbf{Étape 3.} Complète le plan. Le premier jalon est donné en modèle.

\medskip
\textit{Rappel :} un jalon porte une \textbf{date} et un \textbf{état
vérifiable}.

\bigskip

\begin{center}
\begin{tabular}{|c|c|l|}
\hline
\textbf{Jalon} & \textbf{Date} & \textbf{État vérifiable ce jour-là} \\
\hline
J1 & \dotfill & le programme démarre et affiche un menu \\
\hline
J2 & \dotfill & \dotfill \\
\hline
J3 & \dotfill & \dotfill \\
\hline
J4 & \dotfill & \dotfill \\
\hline
\end{tabular}
\end{center}

\bigskip

\textbf{Étape 4.} Ton premier jalon produit-il quelque chose qui \textbf{tourne} ?

\medskip
\begin{center}
\fbox{Oui} \qquad \fbox{Non}
\end{center}

\medskip
\textit{Aide :} si la réponse est non, échange J1 avec un jalon plus simple.

\bigskip

\textbf{Étape 5.} Combien de jours séparent ton dernier jalon de la date de rendu ?

\medskip
\dotfill\ jours

\medskip
\begin{center}
\fbox{c'est une marge suffisante} \qquad \fbox{il faut avancer les jalons}
\end{center}

\bigskip
\bigskip

\subsection*{Partie C --- Se répartir le travail (C2)}

\textbf{Étape 6.} Pour chaque jalon, écris \textbf{un seul} prénom responsable.

\bigskip

\begin{center}
\begin{tabular}{|c|c|}
\hline
\textbf{Jalon} & \textbf{Responsable} \\
\hline
J1 & \dotfill \\
\hline
J2 & \dotfill \\
\hline
J3 & \dotfill \\
\hline
J4 & \dotfill \\
\hline
\end{tabular}
\end{center}

\medskip
\textit{Contrôle :} un jalon sans responsable est un jalon que personne ne fait.
